%%==============================================================================
%% Section III: Preliminaries (notation table, system/threat model, primitives)
%%==============================================================================
\section{Preliminaries}
\label{sec:preliminaries}

\subsection{Notation}
Table~\ref{tab:notation} lists the symbols used throughout the paper. Acronyms are defined at first use, and every mathematical symbol is defined in the table, so subsequent sections do not repeat its definition.

\begin{table}[t]
\centering
\caption{Summary of Key Notation}
\label{tab:notation}
\footnotesize
\setlength{\tabcolsep}{2.5pt}
\renewcommand{\arraystretch}{1.12}
\begin{tabularx}{\columnwidth}{@{}>{\raggedright\arraybackslash}p{0.27\columnwidth} >{\raggedright\arraybackslash}X@{}}
\toprule
\textbf{Notation} & \textbf{Description} \\
\midrule
\multicolumn{2}{@{}l}{\textit{System and Hierarchy}}\\
$\lambda,n,q,\Delta$ & Security parameter, ring dimension, modulus, CKKS scale \\
$N,M,\tau$ & Numbers of worker UAVs, cluster-heads, and threshold \\
$G_c,N_c,c'$ & Worker set of cluster $c$, its size, neighbor head \\
$\ell,T$ & Communication round index and total rounds \\
$S$ & Cloud server \\
\midrule
\multicolumn{2}{@{}l}{\textit{Keys and Shares}}\\
$s_c,b_c$ & CH $c$ secret contribution and public component \\
$s,\bar{b}$ & Group secret and group public key \\
$[s]_c,pk_c^*$ & CH $c$ threshold share and share public key \\
$sk_i^{\mathrm{sig}},vk_i^{\mathrm{sig}}$ & Worker EdDSA signing and verification keys \\
\midrule
\multicolumn{2}{@{}l}{\textit{Updates, Masks, and Hash}}\\
$m_i,m_i'$ & Worker update and blinded update \\
$r_i^{(c)},r_i^{(s)},r_i^{(c')}$ & Three mask shares; $r_i$ is their sum \\
$h_i,\sigma_i$ & Homomorphic hash and worker signature \\
\midrule
\multicolumn{2}{@{}l}{\textit{Ciphertexts, Proofs, and Crypto}}\\
$ct_i',ct^{(c)},ct_{\mathrm{glob}}$ & Individual, cluster, and global ciphertexts \\
$v_i,e_{i,0},e_{i,1}$ & Encryption randomizer and error terms \\
$D_c,\lambda_c$ & Partial decryption and Lagrange coefficient \\
$\pi_i^{\mathrm{enc}},\pi_c$ & Encryption and decryption NIZKs \\
$a,\chi_s,\chi_e$ & Public polynomial, key and error distributions \\
$\mathcal{R}_q,\mathcal{R}_p,\mathrm{negl}(\cdot)$ & Rings and negligible function \\
\bottomrule
\end{tabularx}
\end{table}

\subsection{Hierarchical FL and System Model}
\label{subsec:systemmodel}
We consider a three-tier hierarchical FL system. The $N$ worker UAVs are partitioned into $M$ disjoint clusters $G_1,\ldots,G_M$, with $N_c=|G_c|$ and $\sum_{c=1}^{M}N_c=N$. Each cluster is managed by one cluster-head UAV, and the $M$ heads connect to a cloud server. Worker $i$ holds a local dataset $\mathcal{D}_i$ and a model $w\in\mathbb{R}^{d}$. At round $\ell$, a worker performs local stochastic gradient descent to obtain an update $m_i$; the head pre-aggregates within the cluster, and the cloud aggregates the $M$ cluster results. The correct plaintext aggregate is
\begin{equation}
  \label{eq:plaintext-agg}
  \mglob=\sum_{i=1}^{N} m_i .
\end{equation}
Worker $i$ derives its update through local stochastic gradient descent over $E$ local epochs,
\begin{equation}
  \label{eq:local-update}
  m_i = w_{\ell-1}-w_i^{(E)},\qquad
  w_i^{(k+1)} = w_i^{(k)} - \eta_{\ell}\nabla\mathcal{L}(w_i^{(k)};\mathcal{B}_i),
\end{equation}
where $\mathcal{B}_i$ is a local minibatch and $\mathcal{L}$ the empirical loss. The data are heterogeneous: the class proportion vector $p_k$ of worker $k$ follows $p_k\sim\mathrm{Dir}(\alpha)$, and smaller $\alpha$ yields a more skewed partition.

We use unweighted aggregation for clarity; dataset-size weights $p_i$ can be absorbed by replacing each $m_i$ with $p_i m_i$.

The four role classes are as follows:
\begin{itemize}[leftmargin=*,topsep=2pt,itemsep=1pt]
  \item \textbf{Worker UAV $i$} ($N$, MCU-class, no NTT): local training, one-time pad blinding, sampling of three mask shares, a linear homomorphic hash, and an EdDSA signature.
  \item \textbf{Cluster-head UAV $c$} ($M$, strong): outsourced encryption, NIZK proofs, verifiable secret sharing, intra-cluster pre-aggregation, and partial decryption; it holds one share of the group secret.
  \item \textbf{Neighbor cluster-head $c'$} ($M$, strong): the ring neighbor $c'=(c+1)\bmod M$, which holds the third mask share and has the same capabilities as a head.
  \item \textbf{Cloud server $S$} ($1$, strong): aggregation and verification coordination, including threshold-decryption interpolation; it holds no key share.
\end{itemize}

\subsection{Threat Model and Security Goals}
\label{subsec:threat}
We make the following adversarial assumptions.
\begin{itemize}[leftmargin=*,topsep=2pt,itemsep=1pt]
  \item \textbf{Cloud server.} The cloud is \emph{malicious}: it may fabricate the aggregate, selectively discard updates, tamper with the downlink model, and selectively request the decryption of a chosen individual.
  \item \textbf{Cluster-heads.} Heads are \emph{semi-honest but curious} and corruptible; any coalition of at most $\tau-1$ heads, including a coalition with the cloud, cannot reconstruct the group secret.
  \item \textbf{Worker UAVs.} In the baseline model workers are honest-but-curious and may drop out arbitrarily. In the extended model, a malicious worker may upload any update and a captured worker is fully controlled by the attacker.
  \item \textbf{Mask holders.} Of the three holders of a blinded update, any two that collude cannot recover the update (3-out-of-3).
  \item \textbf{Channels and trust.} No trusted authority exists. The three mask shares are carried over independent secure channels; all other traffic is public.
\end{itemize}
Accordingly, \textsc{VHD-FL} targets four security goals: \emph{G1}, worker privacy against collusion; \emph{G2}, secrecy of the group decryption key against a threshold coalition; \emph{G3}, end-to-end verifiability of encryption, aggregation, and decryption against a malicious cloud; and \emph{G4}, tolerance of arbitrary worker and cluster-head dropouts.

We formalize these goals with two standard games. In the \emph{privacy game}, the adversary chooses two updates $m_0,m_1$, receives the transcript produced for a uniformly hidden bit $b$, and outputs a guess; the protocol is private if its distinguishing advantage is $\negl(\lambda)$. In the \emph{verifiability game}, the adversary controls the cloud and corrupted intermediaries and attempts to make the verifier accept an aggregate different from the honest one; the protocol is verifiable if this success probability is $\negl(\lambda)$. Section~\ref{sec:security} establishes both games and the dropout property.

\subsection{Ring Learning with Errors and CKKS}
\label{subsec:rlwe}
Let $R=\mathbb{Z}[X]/(X^n+1)$ be the cyclotomic ring for a power-of-two $n$, with residue ring $\mathcal{R}_q=R/qR$, and let $\chi_s$ and $\chi_e$ output small polynomials. The \emph{RLWE assumption} states that pairs $(a,b)$ with $a\leftarrow\mathcal{U}(\mathcal{R}_q)$, $s\leftarrow\chi_s$, and $b=[-as+e]_q$ for $e\leftarrow\chi_e$ are computationally indistinguishable from uniform pairs in $\mathcal{R}_q^2$~\cite{lyubashevsky2010ideal}.

CKKS~\cite{cheon2017ckks} maps a real vector to a plaintext polynomial through an isometric isomorphism $\varphi:R\to\mathbb{C}^{n/2}$ and a scaling factor $\Delta$, packing $n/2$ values into one ciphertext (SIMD). It is additively homomorphic, which is all that secure aggregation requires. Recent diagonal packing improves the slot layout for matrix and sparse vectors~\cite{lee2026hydia},~\cite{deryabin2026packing}, and refined modular approximation lowers the noise~\cite{zhang2026efficient}. We use its standard public-key form. The public key is $(\bar b,a)$ with $\bar b=-as+\bar e$; to encrypt a message $m$, sample $v\leftarrow\chi_s$ and $e_0,e_1\leftarrow\chi_e$:
\begin{equation}
  \label{eq:ckks-enc}
  ct=\Enc_{\bar b}(m)=(c_0,c_1)=(v\bar b+m+e_0,\; va+e_1).
\end{equation}
Decryption cancels the $a$-component using the secret:
\begin{equation}
  \label{eq:ckks-dec}
  c_0+s c_1=m+v\bar e+e_0+s e_1\approx m .
\end{equation}
The accumulated small terms are treated as approximation noise. Correct decryption requires this noise to remain below the plaintext resolution:
\begin{equation}
  \label{eq:noise-bound}
  \|v\bar e+e_0+s e_1\|_\infty < \Delta/2 .
\end{equation}
CKKS bootstrapping~\cite{cheon2018bootstrapping} can refresh the noise budget if the depth grows, but our scheme performs only addition, so the noise does not accumulate with depth.

\subsection{Homomorphic Hash}
\label{subsec:hash}
We use a linear homomorphic hash over a group in which finding a non-trivial short relation is SIS-hard. For a vector $m=(m_1,\ldots,m_d)$ and fixed public generators $g_j$,
\begin{equation}
  \label{eq:homhash}
  H(m)=\sum_{j=1}^{d} m_j g_j .
\end{equation}
The map is linear, $H(m+m')=H(m)+H(m')$, and it is collision-resistant: $H(m)=H(m')$ implies $\sum_j(m_j-m_j')g_j=0$, a short non-zero relation among the generators that constitutes a SIS instance~\cite{krohn2004fly},~\cite{micciancio2001svp}. It therefore verifies an aggregate through $H(\mglob)=\sum_i H(m_i)$.

\subsection{Verifiable Secret Sharing and Distributed Key Generation}
\label{subsec:dkg}
A $( \tau, M)$ \emph{Shamir sharing}~\cite{shamir1979} encodes a secret $s$ in a degree-$(\tau-1)$ polynomial $f$ and gives share $f(c)$ to party $c$; any $\tau$ shares in a set $T$ reconstruct $s$ through Lagrange interpolation,
\begin{equation}
  \label{eq:lagrange}
  s=\sum_{c\in T}\lambda_c f(c),\qquad
  \lambda_c=\prod_{j\in T,\,j\neq c}\frac{-j}{c-j},
\end{equation}
while fewer than $\tau$ shares reveal nothing. \emph{Feldman VSS}~\cite{feldman1987} publishes commitments $g^{a_j}$ to the coefficients so that each share is publicly verifiable, and \emph{Pedersen VSS}~\cite{pedersen1991} adds a blinding polynomial that hides the secret information-theoretically. A \emph{distributed key generation} protocol~\cite{gennaro2007dkg} combines these commitments and nonces to produce a joint public key without any dealer, tolerating a dishonest minority. Phase~0 instantiates this protocol among the cluster-heads.

\subsection{Non-Interactive Zero-Knowledge Proofs}
\label{subsec:nizk}
A non-interactive zero-knowledge (NIZK) proof for a relation $\mathcal{R}$ lets a prover demonstrate knowledge of a witness $w$ with $(x,w)\in\mathcal{R}$ without revealing $w$, and it satisfies completeness, knowledge soundness, and zero knowledge. NIZKs are obtained in the random-oracle model through the Fiat--Shamir transform, with succinct pairing-based arguments~\cite{groth2016} or lattice-based instantiations~\cite{grothsahai2008}. We use them to prove, first, that a ciphertext encrypts a committed plaintext and, second, that a partial decryption is correct. All such proofs are generated by the strong cluster-heads, never by the worker UAVs. The only worker-side signature is EdDSA over a twisted Edwards curve, which requires no NTT and completes in tens of microseconds~\cite{rfc8032}.
